\section{Error Analysis and Discussion}

\subsection{Error Analysis}

We categorise errors from the best-performing model (BERT
stage-2) on the MBS test split using keyword-based heuristics to
assign each sentence to a suggestion type. Of 80 test sentences,
BERT stage-2 produces 7 false positives and 6 false negatives (13
errors total).

\paragraph{False positives.} Five of seven false positives fall
  in the imperative category---sentences containing words like
  ``try'', ``recommend'', or ``please''. These are borderline
  cases: our annotation guidelines (following the SemEval
  definition) classify guest-to-guest advice as suggestions, but
  sentences such as ``We would recommend this hotel to everyone''
  are not improvement-oriented from the hotel's perspective. One
  false positive involves a positive modal phrase, and one has no
  keyword signal. These errors highlight the tension between the
  SemEval suggestion definition (any directive) and the practical
  goal of extracting hotel-relevant insights.

\paragraph{False negatives.} Five of six false negatives belong to
  the ``other'' category---suggestions that lack standard keyword
  signals and instead express improvement needs through
  context-dependent phrasing. The remaining false negative is an
  imperative suggestion. The dominance of the ``other'' category
  among false negatives indicates that hotel review suggestions
  often avoid the explicit modal or imperative language that models
  learn to recognise from the SemEval training set.

\subsection{Domain Gap Discussion}

The cross-domain performance drops observed in
Table~\ref{tab:results} stem from systematic linguistic
differences between software forum posts and hotel reviews.
SemEval suggestions tend to be explicit feature requests (``The
system should allow batch uploads''), using modal verbs and
technical vocabulary. Hotel review suggestions are more often
implicit, embedded within narrative praise or complaints (``I wish
the breakfast had more variety''), and use hospitality-specific
language absent from the SemEval training set.

The learned models (TF-IDF~+~LR and BERT stage-1) overfit to
SemEval-specific patterns, resulting in high false-negative rates
on MBS (13 and 12 out of 15 suggestions missed, respectively).
The regex baseline, by contrast, relies on domain-general modal
and conditional patterns that transfer across domains, explaining
its slight improvement on MBS ($+$0.085 suggestion-class F1).

BERT stage-2 partially closes the domain gap. With only 320
in-domain training sentences, it recovers suggestion-class recall
from 0.200 (stage-1) to 0.600 (stage-2), while maintaining
precision at 0.563. This trade-off is visible in the error
comparison: stage-2 reduces false negatives from 12 to 6 but
increases false positives from 3 to 7. In practice, higher recall may be preferable, as
missing genuine suggestions is costlier than reviewing a few
false positives. Figure~\ref{fig:error-comparison} summarises the
FP/FN trade-off across all four models.

\begin{figure}[t]
\centering
\includegraphics[width=\columnwidth]{error_comparison_models.png}
\caption{False positives versus false negatives by model on the
  MBS test split. Stage-2 fine-tuning shifts the error profile
  from FN-dominant to more balanced.}
\label{fig:error-comparison}
\end{figure}

\subsection{Limitations}

\begin{itemize}
  \item \textbf{Annotation scale.} The MBS annotated set contains
  400 sentences, with a test split of only 80 sentences (15
  suggestions). Metrics on this test split have wide confidence
  intervals; a single additional correct prediction would shift
  suggestion-class F1 by approximately 0.05. We did not compute
  bootstrap confidence intervals, so reported differences should
  be treated as indicative rather than statistically confirmed.
  \item \textbf{Single hotel.} All MBS data comes from one
  property. Generalisability to other hotels or hospitality
  segments is unknown.
  \item \textbf{Annotation subjectivity.} Fleiss' $\kappa$ of
  0.379 indicates fair but not strong agreement. The boundary
  between complaints and suggestions is inherently fuzzy, as
  reflected in the 11 calibration ties. Furthermore, 300 of 400
  annotated sentences were labelled by a single annotator each,
  with no inter-annotator validation on these split batches.
  \item \textbf{No human upper bound.} We did not compute
  human--model agreement on the test split.
  \item \textbf{Sentence isolation.} Each sentence is classified
  independently, without review-level context. Suggestions that
  span multiple sentences may be missed.
  \item \textbf{No manual validation.} We did not manually
  validate predictions on the full 70,085-sentence MBS dataset.
  The thematic analysis in Section~5.3 relies on BERT stage-2
  predictions, which achieve suggestion-class F1 of 0.581 on the
  MBS test split but may still contain systematic errors at
  scale.
\end{itemize}

\subsection{Future Directions}

\begin{itemize}
  \item Annotate 1,000+ sentences to reduce test-split variance.
  \item Include multiple hotels for cross-property generalisation.
  \item Investigate few-shot or prompt-based approaches using large
  language models as alternatives to fine-tuning.
  \item Use sentence-pair context to detect implicit suggestions
  that depend on surrounding text.
\end{itemize}
